\documentclass[11pt,a4paper]{article}
\usepackage[top=16mm,bottom=17mm,left=16mm,right=16mm,headheight=0pt,headsep=0mm,footskip=10mm]{geometry}
\usepackage{iftex}
\ifPDFTeX
 \usepackage[utf8]{inputenc}
 \usepackage[T5]{fontenc}
 \usepackage{mathptmx}
\else
 \usepackage{mathptmx}
 \usepackage{fontspec}
 \IfFontExistsTF{Times New Roman}{\setmainfont{Times New Roman}}{\IfFontExistsTF{TeXGyreTermes}{\setmainfont{TeXGyreTermes}}{\setmainfont{Latin Modern Roman}}}
\fi
\usepackage{amsmath,amssymb}
\usepackage{xcolor,tikz,enumitem,multicol,fancyhdr,lastpage,tabularx,array}
\usepackage[most]{tcolorbox}
% Phong cach Toan_11_05.tex; noi dung chep tu de TDM11X17 (15 cau).
\definecolor{questionpurple}{RGB}{137,38,143}
\definecolor{answerblue}{RGB}{35,45,140}
\definecolor{lightblue}{RGB}{238,245,255}
\setlength{\parindent}{0pt}
\setlength{\parskip}{0pt}
\renewcommand{\baselinestretch}{1.10}
\setlength{\columnsep}{8mm}
\pagestyle{fancy}
\fancyhf{}
\renewcommand{\headrulewidth}{0pt}
\fancyfoot[L]{\small Chuyên đề I -- Toán 11}
\fancyfoot[R]{\small Trang \thepage/\pageref{LastPage} -- Bài 06}
\renewcommand{\footrulewidth}{.4pt}
\newtcolorbox{theorybox}{colback=lightblue,colframe=answerblue!60,boxrule=.5pt,
 arc=1mm,left=2.5mm,right=2.5mm,top=1.5mm,bottom=1.5mm,before skip=3pt,after skip=4pt}
\newcommand{\sectitle}[2]{\vspace{4pt}\noindent{\large\bfseries\textcolor{questionpurple}{#1.\ #2}}\par\vspace{3pt}}
% Compact multiple-choice blocks: no blank answer-writing areas from the source PDF.
\newcommand{\question}[3]{\par\addvspace{5pt}\noindent\textcolor{questionpurple}{\textbf{Câu #1.}}\ \textbf{[#2]}\ #3\par\nobreak\vspace{2pt}}
\newcommand{\opt}[2]{\textcolor{answerblue}{\bfseries #1.}\,\ensuremath{#2}}
\newcommand{\choices}[4]{%
 \begingroup\small\renewcommand{\arraystretch}{1.20}%
 \begin{tabularx}{\linewidth}{@{}>{\raggedright\arraybackslash}X@{\hspace{3mm}}>{\raggedright\arraybackslash}X@{}}
 \opt{A}{#1} & \opt{B}{#2}\\
 \opt{C}{#3} & \opt{D}{#4}
 \end{tabularx}\endgroup\par}
\newcommand{\longchoices}[4]{%
 \begingroup\small\par\vspace{1pt}
 \noindent\opt{A}{#1}\par\vspace{2pt}
 \noindent\opt{B}{#2}\par\vspace{2pt}
 \noindent\opt{C}{#3}\par\vspace{2pt}
 \noindent\opt{D}{#4}\par
 \endgroup}
\begin{document}
\noindent
\begin{minipage}[t]{.35\linewidth}
 \centering{\bfseries CHUYÊN ĐỀ I -- TOÁN -- 11\par}
 \vspace{2pt}\rule{.73\linewidth}{.4pt}\par\vspace{3pt}
 CHƯƠNG I\par\vspace{3pt}{\itshape TÀI LIỆU HỌC TẬP -- BÀI 06\par}
\end{minipage}\hfill
\begin{minipage}[t]{.62\linewidth}
 \centering{\bfseries HÀM SỐ LƯỢNG GIÁC\par VÀ PHƯƠNG TRÌNH LƯỢNG GIÁC\par}
 \vspace{4pt}{\small\bfseries BÀI 06. BẤT PHƯƠNG TRÌNH LƯỢNG GIÁC CƠ BẢN\par VÀ TẬP XÁC ĐỊNH CỦA HÀM SỐ LƯỢNG GIÁC\par}
\end{minipage}
\vspace{5mm}
\noindent\textbf{Họ và tên:}\dotfill\hfill\textbf{Lớp:}\enspace\dotfill
\vspace{3mm}
\begin{theorybox}
\textbf{Đề tự luyện TDM11X17:} Bất phương trình lượng giác cơ bản và tập
xác định của hàm số lượng giác cơ bản. \textbf{15 câu trắc nghiệm -- 60 phút.}
Các phương án được giữ theo đề gốc; $k\in\mathbb Z$ xuất hiện trong từng lựa chọn.
\end{theorybox}

\sectitle{A}{Các bài tập nền tảng và vận dụng trọng điểm}

\question{1}{TDM32}{Tất cả các nghiệm của bất phương trình $\sin x\ge\dfrac12$ là}
\choices
 {\dfrac{\pi}{6}+2k\pi\le x\le\dfrac{5\pi}{6}+2k\pi,\ k\in\mathbb Z.}
 {-\dfrac{\pi}{6}+2k\pi\le x\le\dfrac{\pi}{6}+2k\pi,\ k\in\mathbb Z.}
 {\dfrac{\pi}{3}+2k\pi\le x\le\dfrac{2\pi}{3}+2k\pi,\ k\in\mathbb Z.}
 {2k\pi\le x\le\pi+2k\pi,\ k\in\mathbb Z.}

\question{2}{TDM22}{Tất cả các nghiệm của bất phương trình $\sin(2x+3)>2$ là}
\choices
 {\varnothing.}
 {\mathbb R.}
 {x\ge 2k\pi,\ k\in\mathbb Z.}
 {x\ge\arcsin 2+2k\pi,\ k\in\mathbb Z.}

\question{3}{TDM32}{Tập xác định của hàm số $y=\sqrt{-2\sin x-\sqrt2}$ là}
\choices
 {\dfrac{5\pi}{4}+2k\pi\le x\le\dfrac{7\pi}{4}+2k\pi,\ k\in\mathbb Z.}
 {-\dfrac{\pi}{4}+2k\pi\le x\le\dfrac{5\pi}{4}+2k\pi,\ k\in\mathbb Z.}
 {\dfrac{5\pi}{4}+k\pi\le x\le\dfrac{9\pi}{4}+k\pi,\ k\in\mathbb Z.}
 {\dfrac{\pi}{4}+2k\pi\le x\le\dfrac{3\pi}{4}+2k\pi,\ k\in\mathbb Z.}

\question{4}{TDM32}{Tập xác định của hàm số $y=\sqrt{-1-\sin\left(x-\dfrac{\pi}{3}\right)}$ là}
\choices
 {\varnothing.}
 {x=-\dfrac{\pi}{6}+2k\pi,\ k\in\mathbb Z.}
 {\mathbb R.}
 {x=-\dfrac{\pi}{2}+2k\pi,\ k\in\mathbb Z.}

\question{5}{TDM32}{Tập xác định của hàm số $y=\sqrt{2\cos\left(x-\dfrac{\pi}{3}\right)+\sqrt3}$ là}
\choices
 {-\dfrac{5\pi}{6}+2k\pi\le x\le\dfrac{11\pi}{6}+2k\pi,\ k\in\mathbb Z.}
 {-\dfrac{\pi}{2}+2k\pi\le x\le\dfrac{7\pi}{6}+2k\pi,\ k\in\mathbb Z.}
 {\dfrac{5\pi}{6}+2k\pi\le x\le\dfrac{7\pi}{6}+2k\pi,\ k\in\mathbb Z.}
 {-\dfrac{5\pi}{6}+2k\pi\le x\le\dfrac{5\pi}{6}+2k\pi,\ k\in\mathbb Z.}

\question{6}{TDM22}{Tập xác định của hàm số $y=\dfrac1{\sqrt{4-\cos(2025x+2029)}}$ là}
\choices
 {\mathbb R.}
 {\varnothing.}
 {x<\arccos4+2k\pi,\ k\in\mathbb Z.}
 {x>\arccos4+2k\pi,\ k\in\mathbb Z.}

\question{7}{TDM32}{Tập xác định của hàm số $y=\dfrac1{\sqrt{-1-2\cos x}}$ là}
\choices
 {\dfrac{2\pi}{3}+2k\pi<x<\dfrac{4\pi}{3}+2k\pi,\ k\in\mathbb Z.}
 {\dfrac{2\pi}{3}+2k\pi\le x\le\dfrac{4\pi}{3}+2k\pi,\ k\in\mathbb Z.}
 {-\dfrac{2\pi}{3}+2k\pi<x<\dfrac{2\pi}{3}+2k\pi,\ k\in\mathbb Z.}
 {-\dfrac{2\pi}{3}+2k\pi\le x\le\dfrac{2\pi}{3}+2k\pi,\ k\in\mathbb Z.}

\question{8}{TDM32}{Tập xác định của hàm số $y=\sqrt{\cos x-\sqrt3\sin x-2}$ là}
\choices
 {x=\pm\dfrac{\pi}{3}+2k\pi,\ k\in\mathbb Z.}
 {x=-\dfrac{\pi}{3}+2k\pi,\ k\in\mathbb Z.}
 {-\dfrac{\pi}{3}+2k\pi\le x\le\dfrac{\pi}{3}+2k\pi,\ k\in\mathbb Z.}
 {\varnothing.}

\question{9}{TDM32}{Tập xác định của hàm số $y=\dfrac1{\sqrt{\tan\left(2x-\dfrac{\pi}{4}\right)-1}}$ là}
\choices
 {\dfrac{\pi}{2}+k\pi<x<\dfrac{3\pi}{2}+k\pi,\ k\in\mathbb Z.}
 {\dfrac{\pi}{2}+k\pi\le x<\dfrac{3\pi}{2}+k\pi,\ k\in\mathbb Z.}
 {\dfrac{5\pi}{6}+2k\pi\le x\le\dfrac{7\pi}{6}+2k\pi,\ k\in\mathbb Z.}
 {\dfrac{\pi}{4}+\dfrac{k\pi}{2}<x<\dfrac{3\pi}{8}+\dfrac{k\pi}{2},\ k\in\mathbb Z.}

\question{10}{TDM32}{Tập xác định của hàm số $y=\sqrt{\sqrt3-\tan x}$ là}
\choices
 {\dfrac{\pi}{3}+k\pi\le x<\dfrac{\pi}{2}+k\pi,\ k\in\mathbb Z.}
 {-\dfrac{\pi}{2}+k\pi<x\le\dfrac{\pi}{3}+k\pi,\ k\in\mathbb Z.}
 {\mathbb R.}
 {\varnothing.}

\question{11}{TDM32}{Tập xác định của hàm số $y=\sqrt{\cot x+1}$ là}
\choices
 {-\dfrac{\pi}{2}+k\pi<x\le\dfrac{\pi}{4}+k\pi,\ k\in\mathbb Z.}
 {-\dfrac{\pi}{4}+k\pi\le x<\dfrac{\pi}{2}+k\pi,\ k\in\mathbb Z.}
 {k\pi<x\le\dfrac{3\pi}{4}+k\pi,\ k\in\mathbb Z.}
 {\dfrac{3\pi}{4}+k\pi\le x<\pi+k\pi,\ k\in\mathbb Z.}

\question{12}{TDM32}{Tập xác định của hàm số $y=\sqrt{-1-\sqrt3\cot x}$ là}
\choices
 {\dfrac{2\pi}{3}+k\pi\le x<\pi+k\pi,\ k\in\mathbb Z.}
 {k\pi<x\le\dfrac{2\pi}{3}+k\pi,\ k\in\mathbb Z.}
 {\varnothing.}
 {\dfrac{\pi}{3}+k\pi\le x<\pi+k\pi,\ k\in\mathbb Z.}

\question{13}{TDM32}{Tập xác định của hàm số $y=\sqrt{1-2\sin x}\,\sqrt{1-2\cos x}$ là}
\choices
 {-\dfrac{\pi}{3}+2k\pi\le x\le\dfrac{\pi}{6}+2k\pi,\ k\in\mathbb Z.}
 {\dfrac{\pi}{6}+2k\pi\le x\le\dfrac{\pi}{3}+2k\pi,\ k\in\mathbb Z.}
 {\dfrac{\pi}{3}+2k\pi\le x\le\dfrac{5\pi}{6}+2k\pi,\ k\in\mathbb Z.}
 {\dfrac{\pi}{3}+k\pi\le x<\pi+k\pi,\ k\in\mathbb Z.}

% PDF nguon in tu so cau 14 la "sqrt(2 cos + 1)" (khong co bien x).
% Giu nguyen theo anh chup, khong tu y sua thanh "cos x".
\question{14}{TDM32}{Tập xác định của hàm số $y=\dfrac{\sqrt{2\cos+1}}{\sqrt{\sqrt3-2\cos x}}$ là}
\longchoices
 {\dfrac{\pi}{6}+k\pi<x\le\dfrac{2\pi}{3}+k\pi,\ k\in\mathbb Z.}
 {\left[-\dfrac{\pi}{6}+2k\pi;\dfrac{\pi}{6}+2k\pi\right)\cup\left(-\dfrac{2\pi}{3}+2k\pi;\dfrac{2\pi}{6}+2k\pi\right],\ k\in\mathbb Z.}
 {\left(\dfrac{\pi}{6}+2k\pi;\dfrac{2\pi}{3}+2k\pi\right]\cup\left[-\dfrac{2\pi}{3}+2k\pi;-\dfrac{\pi}{6}+2k\pi\right),\ k\in\mathbb Z.}
 {\left[\dfrac{\pi}{6}+2k\pi;\dfrac{2\pi}{3}+2k\pi\right)\cup\left(-\dfrac{2\pi}{3}+2k\pi;-\dfrac{\pi}{6}+2k\pi\right],\ k\in\mathbb Z.}

\question{15}{TDM32}{Tập xác định của hàm số $y=\sqrt{2\sin x-1}\,\sqrt{1-\tan x}$ là}
\longchoices
 {\left[\dfrac{\pi}{4}+2k\pi;\dfrac{\pi}{2}+2k\pi\right],\ k\in\mathbb Z.}
 {\left[\dfrac{\pi}{4}+2k\pi;\dfrac{\pi}{2}+2k\pi\right)\cup\left(\dfrac{\pi}{2}+2k\pi;\dfrac{5\pi}{6}+2k\pi\right],\ k\in\mathbb Z.}
 {\left[\dfrac{\pi}{4}+k\pi;\dfrac{5\pi}{6}+2k\pi\right],\ k\in\mathbb Z.}
 {\left[\dfrac{\pi}{6}+2k\pi;\dfrac{\pi}{4}+2k\pi\right]\cup\left(\dfrac{\pi}{2}+2k\pi;\dfrac{5\pi}{6}+2k\pi\right],\ k\in\mathbb Z.}

\vspace{2mm}
\begin{center}\bfseries\textcolor{answerblue}{--------- Hết ---------}\end{center}
\begin{tcolorbox}[colback=white,colframe=black!35,boxrule=.4pt,arc=0mm,
 left=2mm,right=2mm,top=1mm,bottom=1mm]
\footnotesize
\textbf{Nguồn nội dung:} TƯ DUY MỞ, \emph{Đề tự luyện TDM11X17: Bất phương trình
lượng giác cơ bản và tập xác định của hàm số lượng giác cơ bản},
biên soạn Nguyễn Đăng Ái (PDF gốc 3 trang, 15 câu).
Giữ nguyên nội dung câu hỏi và phương án của bản chụp;
định dạng lại theo \texttt{Toan\_11\_05.tex}, lược bỏ khoảng trống và không kèm lời giải.
\end{tcolorbox}
\end{document}
